\chapter{Návrh riešenia}
\label{ch:solution_approach}

Predmetom tejto kapitoly je koncepčný návrh decentralizovanej aplikácie pre
mikro-tendre študentských rád. Návrh vychádza z~požiadaviek formulovaných
v~kapitole~\ref{ch:problem_statement} a~z~výsledkov analýzy existujúcich riešení
uvedených v~kapitole~\ref{ch:newchapter}. Cieľom nie je navrhnúť univerzálny
systém verejného obstarávania, ale vytvoriť funkčný prototyp, ktorý zodpovedá
špecifickému kontextu malých nákupov študentských organizácií. Navrhované riešenie
preto musí byť technologicky primerané rozsahu bakalárskej práce a~zároveň
adresovať hlavné identifikované problémy: nedostatok transparentnosti, manuálnu
administratívu, slabú auditovateľnosť a~obmedzenú integritu dokumentácie.

Kapitola najprv opisuje celkovú architektúru systému a~zdôvodňuje kľúčové
architektonické rozhodnutia. Následne sa venuje návrhu smart kontraktu, jeho
dátovým štruktúram, rolám a~stavovému modelu. Ďalšie sekcie pokrývajú životný
cyklus tendra, integráciu s~IPFS, návrh používateľského rozhrania a~nakoniec
bezpečnostné a~transparenčné aspekty riešenia.

%%%%%%%%%%%%%%%%%%%%%%%%%%%%%%%%%%%%%%%%%%%%%%%%%%%%%%
\section{Celková architektúra systému}
\label{sec:system_architecture}

%%%%%%%%%%%%%%
\subsection{Prehľad architektúry}
\label{subsec:architecture_overview}

Navrhovaný systém je založený na trojvrstvovej architektúre, ktorá je
v~prostredí Web3 prirodzeným a~praktickým riešením. V~porovnaní s~tradičnými
webovými aplikáciami, kde sa medzi klientom a~databázou typicky nachádza
aplikačný server, tu rozhodujúca časť logiky prebieha priamo v~smart kontrakte
nasadenom na blockchaine. Systém preto neobsahuje samostatný backendový server
ani relačnú databázu.

Architektúru možno rozdeliť na tri základné vrstvy:
\begin{enumerate}
    \item \textbf{Prezentačná vrstva} je tvorená webovou aplikáciou
    implementovanou v~prostredí React. Používateľ komunikuje so systémom
    prostredníctvom grafického rozhrania v~prehliadači, prepája sa s~peňaženkou
    MetaMask a~odosiela transakcie do blockchainovej siete. Súčasťou
    prezentačnej vrstvy je aj celá logika pre zobrazovanie stavov tendrov,
    formulárov, detailov ponúk, správu používateľských rolí a~nahrávanie
    dokumentov do IPFS.

    \item \textbf{Logická vrstva} je reprezentovaná jedným smart kontraktom
    \texttt{MicroTender.sol}, nasadeným na testovacej sieti Polygon Amoy.
    Kontrakt implementuje pravidlá správy tendrov, registráciu dodávateľov,
    ponuky, hlasovanie a~kontrolu oprávnení. Táto vrstva zastupuje funkciu,
    ktorú by v~klasickej webovej architektúre vykonával aplikačný server.
    Na rozdiel od tradičného backendu je však logika smart kontraktu po nasadení
    nemenná, verejne overiteľná a~vykonávaná deterministicky všetkými uzlami
    siete.

    \item \textbf{Dátová vrstva} je rozdelená medzi blockchain a~IPFS.
    Na blockchaine sú uložené kľúčové metadáta potrebné pre auditovateľnosť
    a~vynucovanie pravidiel: identifikátory tendrov, adresy tvorcov a~dodávateľov,
    ceny ponúk, hlasy a~stavové prechody. Väčšie dokumenty a~súbory sú ukladané
    off-chain v~sieti IPFS, pričom na blockchaine sa eviduje iba ich obsahový
    identifikátor CID. Takéto rozdelenie znižuje náklady na gas a~zároveň
    zachováva možnosť overenia integrity dokumentov.
\end{enumerate}

Takto navrhnutá architektúra znižuje počet samostatných komponentov, ktoré je
potrebné vyvíjať a~udržiavať. Systém neobsahuje žiadny middleware, message
broker ani cachovaciu vrstvu. Z~pohľadu bakalárskej práce ide o~vedomý kompromis
medzi technickou jednoduchosťou a~dostatočnou funkcionalitou prototypu.

%%%%%%%%%%%%%%
\subsection{Technologický stack}
\label{subsec:technology_stack}

Pri návrhu riešenia bol zvolený technologický stack, ktorý zodpovedá požiadavke
na jednoduchú, ale funkčnú implementáciu prototypu. Voľba jednotlivých
technológií bola motivovaná najmä ich zrelosťou, kvalitou dokumentácie,
kompatibilitou s~decentralizovaným prostredím a~dostupnosťou pre vývojára
pracujúceho na bakalárskom projekte.

\begin{itemize}
    \item \textbf{React} (verzia 19) slúži na implementáciu prezentačnej
    vrstvy. Jeho komponentový model umožňuje rozdeliť rozhranie na
    znovupoužiteľné časti, čo zjednodušuje údržbu kódu. Pre správu stavu
    aplikácie sa využívajú zabudované React hooks (\texttt{useState},
    \texttt{useEffect}, \texttt{useRef}), bez potreby externého state
    management riešenia.

    \item \textbf{ethers.js} (verzia 6) zabezpečuje komunikáciu medzi
    frontendom a~blockchainom. Knižnica poskytuje rozhranie pre čítanie
    údajov z~kontraktu prostredníctvom view funkcií, podpisovanie
    a~odosielanie transakcií cez pripojenú peňaženku a~spracovanie
    udalostí emitovaných kontraktom. Oproti alternatíve web3.js ponúka
    ethers.js menšiu veľkosť balíka a~modernejšie API.

    \item \textbf{MetaMask} slúži ako kryptografická peňaženka a~nástroj
    na autorizáciu operácií. Každá transakcia vyžaduje explicitný súhlas
    používateľa v~rozhraní MetaMask, čím sa zabezpečuje, že žiadna operácia
    nemôže prebehnúť bez vedomia vlastníka adresy.

    \item \textbf{Solidity} (verzia 0.8.19) je použitá na implementáciu
    smart kontraktu. Táto verzia bola zvolená pre vstavanú ochranu proti
    pretečeniu celočíselných typov (overflow/underflow), ktorú poskytuje
    od verzie 0.8.0, čím odpadá potreba externej knižnice SafeMath.

    \item \textbf{Hardhat} poskytuje vývojové prostredie pre kompiláciu,
    testovanie a~nasadzovanie kontraktu. Jeho výhodou je integrované
    lokálne prostredie Hardhat Network, podpora pre automatizované testy
    v~JavaScripte a~jednoduchá konfigurácia pre viaceré siete.

    \item \textbf{IPFS} a~služba \textbf{Pinata} sú využité na ukladanie
    dokumentácie spojenej s~tendrami. Pinata poskytuje REST API pre
    nahrávanie súborov a~garantuje ich dostupnosť prostredníctvom
    mechanizmu pinning, čím sa predchádza automatickému odstráneniu
    obsahu z~IPFS siete.

    \item \textbf{Tailwind CSS} je využitý na štýlovanie komponentov
    rozhrania. Utility-first prístup umožňuje rýchlu iteráciu dizajnu
    priamo v~JSX kóde bez písania samostatných CSS súborov.

    \item \textbf{Polygon Amoy} je testovacia sieť Polygon blockchainu,
    na ktorej je možné overovať správanie aplikácie bez použitia
    produkčných prostriedkov. V~porovnaní s~hlavnou sieťou Ethereum
    ponúka Polygon výrazne nižšie transakčné poplatky a~rýchlejšie
    potvrdenie blokov, čo ho robí vhodným pre prototypovanie.
\end{itemize}

Uvedený stack pokrýva všetky vrstvy navrhovanej architektúry a~nevyžaduje
prevádzku žiadnej serverovej infraštruktúry okrem samotného blockchainu.

%%%%%%%%%%%%%%
\subsection{Komunikácia medzi komponentmi}
\label{subsec:component_communication}

Komunikácia v~systéme prebieha dvoma hlavnými kanálmi, ktoré zodpovedajú
dvom typom operácií: zápisu na blockchain a~práci s~dokumentmi.

\textbf{Komunikácia s~blockchainom.} Pri operáciách, ktoré menia stav systému
-- napríklad pri vytvorení tendra, podaní ponuky alebo hlasovaní -- frontend
komunikuje cez MetaMask a~knižnicu ethers.js priamo so smart kontraktom. Každá
takáto operácia vyžaduje potvrdenie používateľa v~peňaženke, podpísanie
transakcie privátnym kľúčom a~následné zapísanie do blockchainu. Výsledok
operácie je k~dispozícii po zahrnutí transakcie do bloku, čo v~sieti Polygon
Amoy trvá spravidla niekoľko sekúnd. Smart kontrakt pri každej dôležitej
operácii emituje udalosť (event), ktorú frontend môže spracovať na potvrdenie
úspešnosti transakcie.

\textbf{Čítanie dát z~kontraktu} prebieha prostredníctvom view funkcií, ktoré
nevyžadujú transakciu ani potvrdenie v~peňaženke. Frontend volá tieto funkcie
cez ethers.js a~výsledky sú vrátené priamo z~lokálneho uzla siete. Medzi
najdôležitejšie view funkcie patria \texttt{getTender()},
\texttt{getTenderBids()}, \texttt{getVoteCount()}, \texttt{getUserRole()}
a~\texttt{isRegisteredVendor()}.

\textbf{Komunikácia s~IPFS.} Pri práci s~dokumentmi frontend komunikuje priamo
so službou Pinata prostredníctvom REST API. Súbor je nahraný cez HTTP POST
požiadavku, Pinata vráti identifikátor obsahu CID a~frontend tento CID
následne odovzdá smart kontraktu pri vytváraní tendra. Prístup k~uloženým
dokumentom prebieha cez verejne dostupné IPFS gateway, kde sa CID transformuje
na URL adresu.

\textbf{Režim len na čítanie.} Dôležitým aspektom návrhu je podpora pre
nepripojených používateľov. Ak používateľ nemá pripojenú peňaženku MetaMask,
aplikácia inicializuje read-only RPC provider a~naďalej umožňuje prezeranie
verejných údajov o~tendroch. Používateľ v~tomto režime nemôže vykonávať
transakcie, ale má prístup k~zoznamom tendrov, ich detailom a~ponukám. Táto
vlastnosť podporuje transparentnosť systému, pretože sprístupňuje informácie
aj pozorovateľom, ktorí nechcú alebo nepotrebujú interagovať s~blockchainom.

%%%%%%%%%%%%%%
\subsection{Architektonické rozhodnutia}
\label{subsec:architecture_decisions}

V~procese návrhu boli prijaté viaceré architektonické rozhodnutia, ktoré
zásadne ovplyvňujú výslednú podobu systému. Uvedené rozhodnutia sú výsledkom
posúdenia technických možností, rozsahu bakalárskej práce a~reálnych potrieb
cieľovej skupiny.

\textbf{Absencia backendu.} Systém neobsahuje žiadny serverový backend.
Všetka biznis logika je implementovaná v~smart kontrakte a~frontend komunikuje
s~blockchainom priamo. Tento prístup bol zvolený z~viacerých dôvodov: znižuje
celkovú komplexitu projektu, eliminuje potrebu prevádzky a~zabezpečenia
serverovej infraštruktúry a~zachováva princíp priameho, overiteľného
prístupu k~dátam bez sprostredkovateľa. Pre rozsah prototypu je tento prístup
dostatočný; pri prípadnom nasadení do produkcie by bolo možné doplniť indexačnú
vrstvu pre optimalizáciu čítania dát.

\textbf{Absencia relačnej databázy.} Blockchain slúži ako primárny zdroj
pravdy pre všetky dáta v~systéme. Aktuálny rozsah projektu nevyžaduje off-chain
cache ani fulltextové vyhľadávanie, preto nebola zavedená žiadna databáza.
Toto rozhodnutie zjednodušuje deployment a~eliminuje potrebu synchronizácie
medzi blockchainom a~externým úložiskom.

\textbf{Jeden smart kontrakt namiesto viacerých.} Celá logika systému je
implementovaná v~jednom kontrakte. Táto voľba je podrobnejšie zdôvodnená
v~sekcii~\ref{subsec:monolithic}.

%%%%%%%%%%%%%%%%%%%%%%%%%%%%%%%%%%%%%%%%%%%%%%%%%%%%%%
\section{Návrh smart kontraktu MicroTender}
\label{sec:smart_contracts_design}

Smart kontrakt je rozhodujúcou súčasťou navrhovaného systému. Definuje pravidlá
správy tendrov, riadenie prístupu, životný cyklus obstarávania, logiku
hlasovania aj evidenciu ponúk. Všetky tieto funkcionality sú sústredené
v~jednom kontrakte \texttt{MicroTender.sol}, ktorý je napísaný v~jazyku Solidity
verzie 0.8.19 a~nasadený na sieti Polygon Amoy.

%%%%%%%%%%%%%%
\subsection{Monolitická architektúra kontraktu}
\label{subsec:monolithic}

Celé riešenie je implementované v~jednom smart kontrakte. Alternatívny prístup
by spočíval v~použití viacerých kontraktov alebo factory patternu, kde by každý
tender bol reprezentovaný samostatnou inštanciou kontraktu. Monolitický prístup
bol uprednostnený z~nasledujúcich dôvodov:

\begin{itemize}
    \item \textbf{Nižšia zložitosť.} Všetka podstatná logika sa nachádza na
    jednom mieste, čo uľahčuje porozumenie systému, jeho testovanie
    a~prípadné úpravy.

    \item \textbf{Nižšie náklady na gas.} Pri factory patterne by každé
    vytvorenie tendra vyžadovalo nasadenie nového kontraktu, čo je
    výrazne nákladnejšia operácia než zápis štruktúry do existujúceho
    mapovania.

    \item \textbf{Jednoduchšia správa stavu.} Všetky tendre, ponuky
    a~hlasovanie sú spravované centralizovane v~jednom kontrakte, čo
    eliminuje potrebu koordinácie medzi viacerými kontraktmi a~zjednodušuje
    dotazovanie na globálne štatistiky.

    \item \textbf{Primeranosť rozsahu.} Pre prototyp bakalárskej práce
    zameraný na mikro-tendre s~obmedzeným počtom súčasných tendrov je
    monolitický kontrakt technicky postačujúci.
\end{itemize}

Monolitický kontrakt zároveň lepšie zodpovedá charakteru mikro-tendrov.
Systém pracuje s~menšími obstarávaniami, pri ktorých je dôležitá jednoduchosť,
prehľadnosť a~nízka administratívna réžia. Z~tohto dôvodu nebolo potrebné
implementovať zložitejšie kontraktové väzby, escrow mechanizmy alebo
samostatné spracovanie finančných tokov.

%%%%%%%%%%%%%%
\subsection{Dátové štruktúry}
\label{subsec:data_structures}

Návrh kontraktu je založený na troch hlavných dátových štruktúrach
(\texttt{struct}), ktoré modelujú kľúčové entity systému, a~na súbore
mapovaní (\texttt{mapping}), ktoré zabezpečujú efektívny prístup k~dátam.

\subsubsection{Štruktúra Tender}

Štruktúra \texttt{Tender} reprezentuje samotný tender a~obsahuje nasledujúce
polia:
\begin{itemize}
    \item \texttt{id} (uint256) -- unikátny identifikátor tendra generovaný
    inkrementálnym počítadlom,
    \item \texttt{creator} (address) -- Ethereum adresa člena rady, ktorý
    tender vytvoril,
    \item \texttt{title} (string) -- názov tendra s~validáciou maximálnej
    dĺžky 200 znakov (\texttt{MAX\_STRING\_LENGTH}),
    \item \texttt{description} (string) -- voľný textový opis predmetu tendra,
    \item \texttt{maxBudget} (uint256) -- maximálny rozpočet v~jednotkách wei,
    \item \texttt{category} (string) -- kategória nákupu,
    \item \texttt{ipfsCID} (string) -- obsahový identifikátor dokumentácie
    uloženej v~IPFS; môže byť prázdny, ak tender nemá priloženú
    dokumentáciu,
    \item \texttt{deadline} (uint256) -- Unix timestamp uzávierky prijímania
    ponúk; nastavuje sa pri publikovaní tendra,
    \item \texttt{votingDeadline} (uint256) -- Unix timestamp uzávierky
    hlasovania; nastavuje sa pri spustení hlasovacej fázy,
    \item \texttt{status} (TenderStatus) -- aktuálny stav tendra podľa
    stavového modelu,
    \item \texttt{createdAt} (uint256) -- Unix timestamp vytvorenia tendra.
\end{itemize}

\subsubsection{Štruktúra Bid}

Štruktúra \texttt{Bid} reprezentuje ponuku dodávateľa podanú ku konkrétnemu
tendru:
\begin{itemize}
    \item \texttt{id} (uint256) -- globálny identifikátor ponuky,
    \item \texttt{tenderId} (uint256) -- odkaz na tender, ku ktorému ponuka
    patrí,
    \item \texttt{vendor} (address) -- adresa registrovaného dodávateľa,
    \item \texttt{price} (uint256) -- navrhovaná cena v~wei; kontrakt
    validuje, že cena je kladná a~nepresahuje rozpočet tendra,
    \item \texttt{deliveryTime} (uint256) -- navrhovaná lehota dodania
    v~dňoch,
    \item \texttt{description} (string) -- textový opis ponuky,
    \item \texttt{submittedAt} (uint256) -- Unix timestamp podania ponuky.
\end{itemize}

\subsubsection{Štruktúra VendorApplication}

Štruktúra \texttt{VendorApplication} reprezentuje žiadosť externého subjektu
o~registráciu ako dodávateľ:
\begin{itemize}
    \item \texttt{id} (uint256) -- identifikátor žiadosti,
    \item \texttt{applicant} (address) -- adresa žiadateľa,
    \item \texttt{companyName} (string) -- názov spoločnosti alebo organizácie,
    \item \texttt{contactInfo} (string) -- kontaktné údaje,
    \item \texttt{description} (string) -- opis podnikateľskej činnosti,
    \item \texttt{submittedAt} (uint256) -- čas podania žiadosti,
    \item \texttt{status} (ApplicationStatus) -- stav žiadosti:
    \texttt{Pending}, \texttt{Approved} alebo \texttt{Rejected}.
\end{itemize}

\subsubsection{Mapovania}

Popri štruktúrach kontrakt využíva niekoľko kľúčových mapovaní, ktoré
zabezpečujú efektívne vyhľadávanie a~prístup k~dátam:
\begin{itemize}
    \item \texttt{mapping(uint256 => Tender) public tenders} -- hlavné
    úložisko tendrov indexované podľa ID,
    \item \texttt{mapping(uint256 => Bid[]) public tenderBids} -- pole
    ponúk pre každý tender,
    \item \texttt{mapping(uint256 => mapping(address => bool)) public hasVoted}
    -- evidencia, či daná adresa už hlasovala v~konkrétnom tendri,
    \item \texttt{mapping(uint256 => mapping(uint256 => uint256)) public
    voteCount} -- počet hlasov pre každú ponuku v~rámci tendra,
    \item \texttt{mapping(address => UserRole) public userRoles} -- priradenie
    rolí jednotlivým adresám,
    \item \texttt{mapping(address => bool) public hasRole} -- booleovský
    príznak, či adresa má aktívnu rolu v~systéme,
    \item \texttt{mapping(address => bool) public registeredVendors} --
    evidencia schválených dodávateľov,
    \item \texttt{mapping(uint256 => VendorApplication) public
    vendorApplications} -- žiadosti dodávateľov,
    \item \texttt{mapping(address => uint256) public
    vendorApplicationByAddress} -- vyhľadávanie žiadosti podľa adresy
    žiadateľa.
\end{itemize}

Takto navrhnutý dátový model umožňuje uchovávať na blockchaine len tie údaje,
ktoré sú dôležité pre kontrolu pravidiel a~auditovateľnosť. Väčšie dokumenty
sa neukladajú priamo do kontraktu, ale sú reprezentované iba hodnotou CID,
čím sa výrazne znižujú transakčné náklady.

%%%%%%%%%%%%%%
\subsection{Roly a oprávnenia}
\label{subsec:roles}

Model oprávnení je navrhnutý ako kombinácia formálnych rolí členov rady
a~samostatnej evidencie schválených dodávateľov. Kontrakt nepoužíva externú
knižnicu pre riadenie prístupu (ako napríklad OpenZeppelin AccessControl), ale
implementuje vlastný systém rolí primeraný rozsahu prototypu.

\subsubsection{Administrátor (Admin)}

Administrátor je v~systéme reprezentovaný vlastníkom kontraktu. Pri nasadení
kontraktu je deployer automaticky nastavený ako \texttt{Admin}
a~jeho adresa je uložená v~stavovej premennej \texttt{owner}. Táto rola
má najvyššie oprávnenia: môže prideľovať alebo odoberať roly ostatným
používateľom prostredníctvom funkcií \texttt{grantRole()} a~\texttt{revokeRole()}.
Prístup k~týmto funkciám je chránený modifikátorom \texttt{onlyOwner},
ktorý overuje, že volajúci je zhodný s~uloženou adresou vlastníka.

\subsubsection{Člen rady (Member)}

Člen rady je používateľ, ktorému administrátor explicitne pridelil rolu
prostredníctvom \texttt{grantRole()}. Členovia rady majú oprávnenie vytvárať
a~publikovať tendre, spúšťať hlasovanie, hlasovať o~ponukách, schvaľovať
alebo zamietať žiadosti dodávateľov a~uzatvárať tendre. Tieto oprávnenia
sú vynucované modifikátorom \texttt{onlyMember}, ktorý kontroluje, či
adresa volajúceho má aktívny príznak v~mapovaní \texttt{hasRole} a~jej rola
je minimálne \texttt{Member}.

\subsubsection{Dodávateľ (Vendor)}

Dodávateľ nie je v~návrhu reprezentovaný ako hodnota v~enumerácii
\texttt{UserRole}. Namiesto toho je jeho oprávnenie podávať ponuky výsledkom
schváleného registračného procesu. Dodávateľ podá žiadosť cez funkciu
\texttt{submitVendorApplication()}, člen rady ju schváli prostredníctvom
\texttt{approveVendorApplication()}, a~adresa dodávateľa sa zapíše do
mapovania \texttt{registeredVendors}. Tento prístup bol zvolený zámerne:
oddeľuje interné roly rady od externých účastníkov procesu, čím zjednodušuje
riadenie prístupu. Administrátor alebo samotný dodávateľ môže status
kedykoľvek odvolať prostredníctvom funkcie \texttt{revokeVendorStatus()},
pričom po odvolaní má dodávateľ možnosť podať novú žiadosť.

\subsubsection{Pozorovateľ (Observer)}

Verejnosť v~kontrakte nemá explicitnú rolu. Ide o~ľubovoľného používateľa,
ktorý pristupuje k~verejne čitateľným dátam kontraktu bez toho, aby disponoval
akoukoľvek priradenou rolou. V~navrhovanom systéme takýto používateľ
môže prezerať zoznamy tendrov, ich detaily a~ponuky prostredníctvom view
funkcií kontraktu. Frontend podporuje tento režim aj bez pripojenej peňaženky,
prostredníctvom read-only RPC providera.

Takto navrhnutý model oprávnení je pre potreby mikro-tendrov študentskej rady
dostatočný. Rola administrátora zodpovedá predsedovi alebo správcovi rady,
členovia rady sú osoby s~hlasovacím právom a~dodávatelia sú externé subjekty
uchádzajúce sa o~zákazku.

%%%%%%%%%%%%%%
\subsection{Stavový model tendra}
\label{subsec:state_model}

Tender prechádza v~systéme viacerými stavmi, ktoré sú definované enumeráciou
\texttt{TenderStatus}. Navrhnutý stavový model obsahuje šesť stavov:
\begin{itemize}
    \item \texttt{Draft} -- tender bol vytvorený, ale ešte nie je verejne
    publikovaný. V~tomto stave je možné aktualizovať CID dokumentácie cez
    \texttt{updateTenderIPFSCID()},
    \item \texttt{Open} -- tender je zverejnený a~aktívne prijíma ponuky od
    registrovaných dodávateľov. Lehota na predkladanie ponúk je definovaná
    parametrom \texttt{deadline},
    \item \texttt{Voting} -- prijímanie ponúk bolo ukončené a~prebieha
    hlasovanie členov rady. Lehota hlasovania je určená parametrom
    \texttt{votingDeadline},
    \item \texttt{Completed} -- hlasovanie bolo ukončené a~víťazná ponuka
    bola určená na základe najvyššieho počtu hlasov,
    \item \texttt{Fulfilled} -- tender bol administratívne uzavretý
    tvorcom ako splnený,
    \item \texttt{Cancelled} -- tender bol zrušený tvorcom. Zrušenie je
    možné iba v~stavoch \texttt{Draft} a~\texttt{Open}.
\end{itemize}

Prechody medzi stavmi sú viazané na konkrétne funkcie kontraktu. Funkcia
\texttt{createTender()} vytvorí tender v~stave \texttt{Draft}. Funkcia
\texttt{publishTender()} presunie tender do stavu \texttt{Open} a~nastaví
uzávierku, pričom parameter \texttt{\_daysUntilDeadline} musí byť
v~rozsahu definovanom konštantami \texttt{MIN\_DEADLINE\_DAYS} (3 dni)
a~\texttt{MAX\_DEADLINE\_DAYS} (30 dní). Funkcia \texttt{startVoting()}
presunie tender do stavu \texttt{Voting} a~nastaví hlasovaciu lehotu
podľa parametra \texttt{\_votingDays}, ktorý musí byť v~rozsahu
\texttt{MIN\_VOTING\_DAYS} (1 deň) až \texttt{MAX\_VOTING\_DAYS}
(14 dní). Po uplynutí hlasovacej lehoty sa tender finalizuje cez
\texttt{finalizeTender()} a~prejde do stavu \texttt{Completed}. Následne
ho tvorca môže označiť ako splnený funkciou \texttt{fulfillTender()}.
V~stavoch \texttt{Draft} a~\texttt{Open} je zároveň možné tender zrušiť
cez \texttt{cancelTender()}.

Každá funkcia meniaca stav tendra overuje platnosť prechodu
prostredníctvom \texttt{require} podmienok. Tým sa zabezpečuje, že tender
nemôže preskočiť žiadnu fázu procesu ani sa vrátiť do predchádzajúceho
stavu. Tento mechanizmus je kľúčový pre zachovanie konzistencie
a~predvídateľnosti obstarávacieho procesu.

%%%%%%%%%%%%%%%%%%%%%%%%%%%%%%%%%%%%%%%%%%%%%%%%%%%%%%
\section{Životný cyklus tendra a proces obstarávania}
\label{sec:tender_lifecycle}

Na úrovni používateľského procesu možno návrh systému rozdeliť do troch
hlavných etáp: vytvorenie a~publikovanie tendra, registrácia dodávateľa
a~podanie ponuky, hlasovanie a~ukončenie tendra. Každá etapa zahŕňa
interakciu medzi jedným alebo viacerými aktérmi a~smart kontraktom.

%%%%%%%%%%%%%%
\subsection{Etapa 1: Vytvorenie a publikovanie tendra}
\label{subsec:lifecycle_creation}

Proces začína členom rady, ktorý identifikuje potrebu obstarávania
a~pripraví základné údaje o~tendri. Vo webovej aplikácii vyplní formulár
obsahujúci názov, textový opis, maximálny rozpočet, kategóriu
a~voliteľne priloží dokumentáciu.

Ak je k~tendru priložená dokumentácia, používateľ ju najprv nahrá do IPFS
prostredníctvom služby Pinata. Frontend zabezpečí validáciu typu a~veľkosti
súboru (maximálne 10~MB, podporované formáty PDF, DOC a~DOCX). Po úspešnom
nahratí frontend obdrží CID, ktorý sa uloží do formulára a~neskôr odovzdá
kontraktu.

V~návrhu aplikácie je tento proces riešený dvojkrokovo. Najprv sa tender
vytvorí ako koncept pomocou funkcie \texttt{createTender()}, čím sa uloží
v~stave \texttt{Draft}. Táto funkcia prijíma päť parametrov: názov, opis,
rozpočet v~wei, kategóriu a~CID dokumentácie. Kontrakt validuje, že názov
nie je prázdny a~nepresahuje \texttt{MAX\_STRING\_LENGTH} znakov, rozpočet
je kladný a~kategória je zadaná. Po úspešnom vytvorení sa emituje udalosť
\texttt{TenderCreated} s~identifikátorom tendra, adresou tvorcu a~názvom.

V~druhom kroku tvorca tender publikuje pomocou \texttt{publishTender()},
pri ktorom zadá lehotu na predkladanie ponúk v~dňoch. Stav sa zmení na
\texttt{Open} a~emituje sa udalosť \texttt{TenderPublished}. Takéto
rozdelenie umožňuje tvorcovi ešte pred publikovaním skontrolovať správnosť
údajov, prípadne aktualizovať CID dokumentácie prostredníctvom funkcie
\texttt{updateTenderIPFSCID()}, ktorá je dostupná výlučne v~stave
\texttt{Draft}.

%%%%%%%%%%%%%%
\subsection{Etapa 2: Registrácia dodávateľa a podanie ponuky}
\label{subsec:lifecycle_bidding}

Dodávateľ sa do systému nezapája automaticky pripojením peňaženky, ale musí
najprv podať žiadosť o~registráciu. Táto kontrola je v~návrhu dôležitá,
pretože umožňuje študentskej rade overovať, kto je oprávnený vstupovať
do procesu obstarávania.

Dodávateľ vyplní formulár s~názvom spoločnosti, kontaktnými údajmi
a~opisom činnosti, a~odošle žiadosť prostredníctvom funkcie
\texttt{submitVendorApplication()}. Kontrakt pritom overuje, že žiadateľ
ešte nie je registrovaným dodávateľom a~nemá rozpracovanú žiadosť.
Po odoslaní sa emituje udalosť \texttt{VendorApplicationSubmitted}.

Člen rady následne žiadosť posudzuje a~rozhodne o~jej schválení alebo
zamietnutí. Schválením cez \texttt{approveVendorApplication()} sa adresa
žiadateľa zapíše do mapovania \texttt{registeredVendors} a~emitujú sa
udalosti \texttt{VendorApplicationApproved} a~\texttt{VendorRegistered}.
Zamietnutím cez \texttt{rejectVendorApplication()} sa stav žiadosti zmení
na \texttt{Rejected} a~žiadateľovi sa umožní podať novú žiadosť, pretože
sa resetuje jeho záznam v~mapovaní \texttt{vendorApplicationByAddress}.

Po schválení registrácie môže dodávateľ podávať ponuky k~aktívnym tendrom
pomocou funkcie \texttt{submitBid()}. Pri podávaní ponuky zadáva cenu v~wei,
lehotu dodania v~dňoch a~textový opis. Kontrakt validuje, že tender je
v~stave \texttt{Open}, termín uzávierky ešte neuplynul, cena je kladná
a~nepresahuje maximálny rozpočet tendra, lehota dodania je kladná
a~volajúci je registrovaný dodávateľ. Po úspešnom podaní sa emituje
udalosť \texttt{BidSubmitted}.

V~aktuálnom prototype sa dokumentácia viaže primárne na samotný tender;
samostatné pole CID pre každú ponuku nie je v~štruktúre \texttt{Bid}
implementované. Toto je vedomé zjednodušenie, ktoré by bolo možné rozšíriť
v~budúcich verziách.

%%%%%%%%%%%%%%
\subsection{Etapa 3: Hlasovanie a ukončenie tendra}
\label{subsec:lifecycle_voting}

Po skončení fázy prijímania ponúk môže tvorca tendra spustiť hlasovanie
pomocou funkcie \texttt{startVoting()}. Podmienkou spustenia je, že tender
je v~stave \texttt{Open} a~obsahuje aspoň jednu podanú ponuku. Tvorca
zároveň zadáva počet dní na hlasovanie, pričom hodnota musí byť v~rozsahu
1 až 14 dní podľa konštánt kontraktu. Tým sa tender presunie do stavu
\texttt{Voting} a~nastaví sa \texttt{votingDeadline}.

V~hlasovacej fáze majú členovia rady prístup k~zoznamu predložených ponúk.
Každý oprávnený člen môže odovzdať hlas práve jednej ponuke prostredníctvom
funkcie \texttt{castVote()}. Kontrakt kontroluje, či tender je v~stave
\texttt{Voting}, hlasovacie obdobie ešte neuplynulo, volajúci má rolu člena
rady, dosiaľ nehlasoval v~tomto tendri a~zvolená ponuka skutočne existuje
v~poli \texttt{tenderBids}. Po úspešnom hlasovaní sa emituje udalosť
\texttt{VoteCasted}.

Overenie existencie ponuky prebieha iteráciou cez pole ponúk tendra
s~optimalizovaným cyklom, kde sa inkrementácia indexu vykonáva v~bloku
\texttt{unchecked} pre zníženie spotreby gas.

Po uplynutí hlasovacej lehoty sa tender ukončí funkciou
\texttt{finalizeTender()}. Kontrakt v~tomto kroku iteruje cez všetky ponuky,
porovnáva počty hlasov a~určí víťaznú ponuku na základe najvyššieho počtu
hlasov. Podmienkou úspešnej finalizácie je, že aspoň jedna ponuka obdržala
aspoň jeden hlas. Po finalizácii sa tender presunie do stavu
\texttt{Completed} a~emituje sa udalosť \texttt{TenderCompleted}
s~identifikátorom víťaznej ponuky. Výsledky hlasovania je možné následne
overiť prostredníctvom funkcie \texttt{getWinningBid()}, ktorá vráti
identifikátor víťaznej ponuky, adresu dodávateľa, cenu a~počet hlasov.

Tvorca tendra môže následne administratívne označiť tender ako splnený
prostredníctvom \texttt{fulfillTender()}, čím sa stav zmení na
\texttt{Fulfilled}.

Je potrebné zdôrazniť, že aktuálna verzia riešenia nepokrýva automatizované
on-chain platby, escrow mechanizmus ani riešenie sporov. Samotné plnenie
zákazky a~finančné vyrovnanie prebiehajú mimo blockchain, čo zodpovedá
reálnej praxi mikro-tendrov v~študentských organizáciách. Návrh sa sústreďuje
na transparentnú evidenciu, správu ponúk a~demokratické hlasovanie, čo je
pre rozsah bakalárskej práce primerané.

%%%%%%%%%%%%%%%%%%%%%%%%%%%%%%%%%%%%%%%%%%%%%%%%%%%%%%
\section{Návrh integrácie s IPFS}
\label{sec:ipfs_design}

Pre ukladanie dokumentácie k~tendrom bol zvolený systém IPFS (InterPlanetary
File System). Hlavným dôvodom je nevhodnosť priameho ukladania väčších
súborov do blockchainu, kde náklady na úložisko rastú lineárne s~veľkosťou
dát. Kým blockchain je určený na nemenné uloženie kľúčových metadát
a~logiky procesu, IPFS je vhodnejšie pre samotné dokumenty, prílohy
a~ďalšie rozsiahlejšie súbory.

%%%%%%%%%%%%%%
\subsection{Dôvod použitia IPFS}
\label{subsec:ipfs_rationale}

Použitie IPFS prináša v~navrhovanom riešení viacero výhod:
\begin{itemize}
    \item \textbf{Zníženie nákladov.} Dokumenty sa neukladajú priamo do
    blockchainu. Na reťazci sa uchováva iba krátky reťazec CID, zatiaľ
    čo samotný súbor môže mať veľkosť niekoľko megabajtov.

    \item \textbf{Overenie integrity.} Každý súbor v~IPFS je
    identifikovaný hodnotou CID, ktorá je odvodená z~obsahu súboru
    prostredníctvom kryptografickej hašovacej funkcie. Akákoľvek zmena
    obsahu dokumentu vedie k~odlišnému CID. Blockchain preto uchováva
    kryptografický dôkaz o~tom, aký presný obsah bol k~tendru pripojený
    v~čase jeho vytvorenia.

    \item \textbf{Decentralizácia úložiska.} IPFS je distribuovaná
    sieť, kde súbory môžu byť replikované na viacerých uzloch. Aj keď
    aktuálny prototyp závisí od služby Pinata pre garantovanú
    dostupnosť, samotný protokol umožňuje prístup k~súboru z~ľubovoľného
    uzla, ktorý ho uchováva.

    \item \textbf{Prepojenie s~blockchainom.} Uloženie CID na blockchaine
    vytvára auditovateľnú väzbu medzi tendrom a~jeho dokumentáciou.
    Ktokoľvek môže stiahnuť dokument z~IPFS a~overiť, že jeho CID
    zodpovedá hodnote uloženej v~kontrakte.
\end{itemize}

%%%%%%%%%%%%%%
\subsection{Proces nahrávania cez Pinata}
\label{subsec:pinata_upload}

V~navrhovanom riešení prebieha nahrávanie dokumentu priamo z~frontendu
cez službu Pinata, bez potreby samostatného backendového servera.
Proces pozostáva z~nasledujúcich krokov:
\begin{enumerate}
    \item Používateľ vo formulári vyberie súbor, ktorý chce priložiť
    k~tendru.
    \item Frontend vykoná validáciu: overí, že veľkosť súboru nepresahuje
    10~MB a~typ súboru zodpovedá povoleným formátom (PDF, DOC, DOCX).
    \item Súbor sa odošle na Pinata API prostredníctvom HTTP POST
    požiadavky na endpoint \texttt{/pinning/pinFileToIPFS}. Autentifikácia
    prebieha prostredníctvom JWT tokenu uloženého v~premennej prostredia.
    \item Pinata súbor spracuje, uloží do IPFS siete a~vráti CID.
    \item Frontend obdržaný CID uchová a~pri vytváraní tendra ho odovzdá
    smart kontraktu ako parameter funkcie \texttt{createTender()}.
\end{enumerate}

Takto navrhnutý postup má dve dôležité vlastnosti. Po prvé, nie je
potrebné prevádzkovať samostatný backend len na účely uploadu súborov --
celá komunikácia s~IPFS prebieha na strane klienta. Po druhé, používateľ
vníma celý proces ako súčasť jednej aplikácie, hoci technicky ide
o~kombináciu komunikácie s~IPFS a~blockchainom.

%%%%%%%%%%%%%%
\subsection{Prístup k dokumentom}
\label{subsec:document_access}

Po uložení CID do kontraktu je možné dokument získať späť prostredníctvom
IPFS gateway. Frontend preto nepotrebuje uchovávať vlastnú kópiu dokumentu
ani ho ukladať do ďalšieho úložiska. V~detaile tendra môže používateľ
dokument otvoriť alebo stiahnuť prostredníctvom URL adresy odvodenej z~CID,
napríklad vo formáte \texttt{https://ipfs.io/ipfs/\{CID\}}.

Aplikácia používa utilitu \texttt{getIPFSUrl()}, ktorá transformuje CID
na prístupovú URL adresu. Tento prístup je transparentný a~univerzálny --
akýkoľvek používateľ s~prístupom k~CID môže dokument overiť nezávisle
na frontende aplikácie, prostredníctvom ľubovoľného IPFS klienta alebo
verejného gateway.

%%%%%%%%%%%%%%%%%%%%%%%%%%%%%%%%%%%%%%%%%%%%%%%%%%%%%%
\section{Návrh používateľského rozhrania (UI/UX)}
\label{sec:ui_design}

Používateľské rozhranie bolo navrhnuté s~dôrazom na to, aby skrylo
väčšinu technickej zložitosti blockchainu a~IPFS. Hoci systém využíva
decentralizované technológie, používateľ by mal mať pocit práce s~bežnou
modernou webovou aplikáciou. Návrh vychádza z~predpokladu, že primárnou
cieľovou skupinou sú členovia študentskej rady a~lokálni dodávatelia, ktorí
nemusia mať predchádzajúce skúsenosti s~kryptomenami ani blockchainovou
infraštruktúrou.

%%%%%%%%%%%%%%
\subsection{Hlavné obrazovky aplikácie}
\label{subsec:screens}

Návrh rozhrania vychádza z~konkrétnych obrazoviek, z~ktorých každá zodpovedá
jednej funkčnej oblasti systému:
\begin{itemize}
    \item \textbf{Dashboard} poskytuje celkový prehľad o~tendroch a~stave
    systému. Zobrazuje súhrnné štatistiky, zoznamy najnovších tendrov
    a~rýchle odkazy na najčastejšie používané funkcie podľa roly
    prihláseného používateľa.

    \item \textbf{CreateTender} je formulárová obrazovka, kde člen rady
    vytvára nový tender. Zahŕňa dvojkrokový workflow: v~prvom kroku sa
    vyplnia údaje a~prípadne nahrá dokumentácia, v~druhom kroku sa tender
    publikuje s~definovanou lehotou.

    \item \textbf{MyTenders} zobrazuje zoznam tendrov vytvorených
    prihláseným členom rady, čím umožňuje rýchly prístup k~vlastným
    tendrom a~ich správu.

    \item \textbf{AllTenders} zobrazuje kompletný zoznam všetkých tendrov
    v~systéme bez ohľadu na tvorcu. Táto obrazovka je dostupná aj
    nepripojeným používateľom v~režime len na čítanie.

    \item \textbf{TenderDetail} zobrazuje detail konkrétneho tendra
    vrátane jeho metadát, priloženej dokumentácie, zoznamu ponúk
    a~možnosti hlasovania alebo podávania ponúk podľa roly a~stavu
    tendra.

    \item \textbf{Voting} je špecializovaná obrazovka pre hlasovanie
    členov rady, kde sú prehľadne zobrazené tendre vo fáze hlasovania
    spolu s~predloženými ponukami.

    \item \textbf{VendorRegistration} umožňuje externému subjektu podať
    žiadosť o~registráciu ako dodávateľ. Obrazovka zobrazuje aj aktuálny
    stav prípadnej existujúcej žiadosti.

    \item \textbf{VendorApproval} je určená pre členov rady a~slúži na
    posudzovanie žiadostí dodávateľov -- schvaľovanie, zamietanie
    a~prípadné odvolanie statusu dodávateľa.
\end{itemize}

Navigácia v~systéme je riadená cez komponenty \texttt{Sidebar}
a~\texttt{Header}. \texttt{Sidebar} zobrazuje položky navigácie
filtrované podľa roly prihláseného používateľa, \texttt{Header}
poskytuje informácie o~pripojenom účte a~mobilnú navigáciu.

%%%%%%%%%%%%%%
\subsection{Adaptívne rozhranie podľa roly používateľa}
\label{subsec:adaptive_ui}

Jedným z~dôležitých prvkov návrhu je prispôsobenie rozhrania rôznym typom
používateľov. Nie všetky obrazovky a~operácie sú dostupné každému. Systém
rozlišuje nasledujúce úrovne prístupu:

\begin{itemize}
    \item \textbf{Člen rady} má plný prístup k~vytváraniu tendrov,
    hlasovaniu, schvaľovaniu žiadostí dodávateľov a~správe vlastných
    tendrov. V~navigácii vidí položky Nový tender, Moje tendery,
    Hlasovanie a~Žiadosti dodávateľov.

    \item \textbf{Registrovaný dodávateľ} má prístup k~prezeraniu tendrov,
    ich detailom a~podávaniu ponúk. V~detaile tendra v~stave \texttt{Open}
    vidí formulár na odoslanie ponuky.

    \item \textbf{Neregistrovaný používateľ s~peňaženkou} vidí možnosť
    podať žiadosť o~registráciu ako dodávateľ prostredníctvom obrazovky
    VendorRegistration.

    \item \textbf{Nepripojený používateľ} (bez peňaženky) môže systém
    využívať v~režime len na čítanie. Vidí Dashboard, AllTenders a~detaily
    tendrov, ale nemôže vykonávať žiadne transakcie.
\end{itemize}

Toto prispôsobenie je implementované na úrovni frontendu, kde sa na
základe načítanej roly z~kontraktu podmienene zobrazujú alebo skrývajú
jednotlivé obrazovky a~ovládacie prvky. Ak používateľ nemá oprávnenie
na prístup ku konkrétnej obrazovke, je automaticky presmerovaný na
Dashboard.

%%%%%%%%%%%%%%
\subsection{UX princípy}
\label{subsec:ux_principles}

Pri návrhu rozhrania boli zohľadnené nasledujúce zásady, ktoré sú osobitne
dôležité v~kontexte blockchain aplikácií:

\begin{itemize}
    \item \textbf{Minimalizácia blockchainovej zložitosti.} Používateľ
    nemusí rozumieť detailom práce s~kontraktom, hašmi, wei jednotkami
    alebo RPC komunikáciou. Ceny sú zobrazované v~EUR s~automatickým
    prepočtom na wei, adresy sú skrátené na čitateľný formát a~stavy
    tendrov sú prezentované pomocou zrozumiteľných slovenských názvov
    (Koncept, Aktívny, Hlasovanie, Ukončený, Splnený, Zrušený).

    \item \textbf{Dvojkrokové workflow.} Vytvorenie tendra je rozdelené
    na dva kroky: draft a~publikovanie. Tento návrh znižuje
    pravdepodobnosť chýb, pretože tvorca môže skontrolovať údaje pred
    ich konečným zverejnením a~prípadne upraviť priloženú dokumentáciu.

    \item \textbf{Spätná väzba pri operáciách.} Pri každej transakcii
    sa zobrazuje indikátor načítavania, ktorý informuje používateľa
    o~tom, že prebieha spracovanie a~čaká sa na potvrdenie blockchainom.
    Po dokončení operácie je používateľ informovaný výsledkom
    prostredníctvom dialógového okna.

    \item \textbf{Zrozumiteľnosť stavov.} Systém zobrazuje aktuálny stav
    tendra, žiadosti dodávateľa a~výsledky hlasovania tak, aby bolo
    zrejmé, v~akej fáze sa proces nachádza. Každý stav je vizuálne
    odlíšený farebným označením.

    \item \textbf{Validácia vstupov.} Formuláre vykonávajú validáciu
    vstupných údajov ešte pred odoslaním transakcie, čím sa predchádza
    zbytočným pokusom o~transakciu, ktoré by boli kontraktom zamietnuté.

    \item \textbf{Automatické prepínanie siete.} Pri pripájaní peňaženky
    systém overí, či je používateľ pripojený k~správnej sieti (Polygon
    Amoy). Ak nie, automaticky ponúkne prepnutie alebo pridanie siete
    do MetaMask.
\end{itemize}

%%%%%%%%%%%%%%%%%%%%%%%%%%%%%%%%%%%%%%%%%%%%%%%%%%%%%%
\section{Bezpečnosť a transparentnosť}
\label{sec:security}

Navrhované riešenie je postavené na princípe, že dôležité pravidlá
obstarávania majú byť vynucované technicky, nie len organizačne.
Bezpečnosť a~transparentnosť preto netvoria samostatnú doplnkovú
funkciu systému, ale sú súčasťou jeho základného návrhu.

%%%%%%%%%%%%%%
\subsection{Bezpečnostné vlastnosti}
\label{subsec:security_properties}

Z~pohľadu bezpečnosti sú v~návrhu zohľadnené nasledujúce vlastnosti:

\begin{itemize}
    \item \textbf{Kontrola prístupu.} Prístup ku kritickým funkciám je
    riadený cez modifikátory \texttt{onlyOwner} a~\texttt{onlyMember}.
    Modifikátor \texttt{onlyOwner} obmedzuje správu rolí na vlastníka
    kontraktu, modifikátor \texttt{onlyMember} obmedzuje operácie
    spojené s~tendrami na členov rady. Dodávateľ môže podať ponuku
    len po schválení registrácie, čo je overované podmienkou
    \texttt{registeredVendors[msg.sender]}.

    \item \textbf{Validácia vstupov.} Každá verejná funkcia kontraktu
    obsahuje \texttt{require} podmienky, ktoré validujú vstupné údaje.
    Overuje sa platnosť identifikátorov cez modifikátory
    \texttt{validTender} a~\texttt{validApplication}, správnosť stavov,
    dodržanie lehôt, nenulovosť adries a~maximálna dĺžka reťazcov
    (\texttt{MAX\_STRING\_LENGTH}).

    \item \textbf{Ochrana proti pretečeniu.} Solidity vo verzii 0.8.19
    poskytuje vstavanú kontrolu pretečenia a~podtečenia celočíselných
    typov. V~kontrakte sú zároveň využité bloky \texttt{unchecked} na
    miestach, kde je pretečenie logicky vylúčené (inkrementácia
    počítadiel), čo znižuje spotrebu gas bez vplyvu na bezpečnosť.

    \item \textbf{Prevencia duplicitného hlasovania.} Mapovanie
    \texttt{hasVoted[\_tenderId][msg.sender]} zabraňuje tomu, aby
    jeden člen hlasoval v~tom istom tendri viackrát.

    \item \textbf{Prevencia duplicitnej registrácie.} Kontrakt overuje,
    že žiadateľ ešte nie je registrovaným dodávateľom a~nemá
    rozpracovanú žiadosť, čím sa zabráni opakovanému podávaniu
    žiadostí.

    \item \textbf{Obmedzenie na tvorcu tendra.} Funkcie
    \texttt{publishTender()}, \texttt{cancelTender()},
    \texttt{startVoting()}, \texttt{finalizeTender()} a~\texttt{fulfillTender()}
    sú obmedzené na adresu, ktorá tender vytvorila, prostredníctvom
    podmienky \texttt{tender.creator == msg.sender}.
\end{itemize}

Návrh zároveň zámerne neobsahuje zložitejšie bezpečnostné mechanizmy,
ako sú anonymné hlasovanie, on-chain escrow, multisig správa alebo
automatizované platby. Ich neprítomnosť nie je chybou návrhu, ale vedomým
zúžením rozsahu na jadrové vlastnosti systému, ktoré je možné v~rámci
bakalárskej práce implementovať a~otestovať spoľahlivo. Tieto rozšírenia
sú diskutované v~záverečnej kapitole ako možný smer budúcej práce.

%%%%%%%%%%%%%%
\subsection{Transparentnosť a auditovateľnosť}
\label{subsec:transparency}

Transparentnosť systému je zabezpečená tým, že kľúčové operácie sú
zapisované do blockchainu a~zanechávajú auditovateľnú stopu. Kontrakt
emituje nasledujúce udalosti (events):
\begin{itemize}
    \item \texttt{TenderCreated} -- pri vytvorení tendra,
    \item \texttt{TenderPublished} -- pri publikovaní tendra,
    \item \texttt{TenderCancelled} -- pri zrušení tendra,
    \item \texttt{BidSubmitted} -- pri podaní ponuky,
    \item \texttt{VoteCasted} -- pri odovzdaní hlasu,
    \item \texttt{TenderCompleted} -- pri finalizácii tendra,
    \item \texttt{VendorApplicationSubmitted} -- pri podaní žiadosti
    dodávateľa,
    \item \texttt{VendorApplicationApproved} -- pri schválení žiadosti,
    \item \texttt{VendorApplicationRejected} -- pri zamietnutí žiadosti,
    \item \texttt{RoleGranted} a~\texttt{RoleRevoked} -- pri zmene rolí,
    \item \texttt{IPFSCIDUpdated} -- pri aktualizácii dokumentácie tendra.
\end{itemize}

Tieto udalosti sú uložené v~transakčných logoch blockchainu a~sú
dostupné prostredníctvom blockchain explorerov, napríklad Polygonscan
Amoy. Ľubovoľný používateľ tak môže spätne overiť, kto tender vytvoril,
kedy bol publikovaný, aké ponuky boli podané, kto hlasoval a~ktorá
ponuka zvíťazila.

Dôležitú úlohu v~transparentnosti zohráva aj prepojenie blockchainu
a~IPFS. Dokumentácia tendra nie je len pripojená ako externá príloha,
ale je reprezentovaná identifikátorom CID, ktorý umožňuje overiť, či
sa obsah dokumentu od času uloženia nezmenil. Tým sa posilňuje
dôveryhodnosť celého obstarávacieho procesu -- zmena dokumentu by viedla
k~odlišnému CID, čo by bolo okamžite zistiteľné porovnaním s~hodnotou
uloženou na blockchaine.

Výsledkom návrhu je riešenie, ktoré nie je zamerané na úplnú automatizáciu
všetkých ekonomických a~administratívnych aspektov verejného obstarávania,
ale na transparentnú a~technicky konzistentnú podporu mikro-tendrov
v~prostredí študentskej rady. Všetky dôležité udalosti sú zaznamenané
nemenným spôsobom, pravidlá sú vynucované kódom kontraktu a~dokumentácia
je prepojená s~blockchainom prostredníctvom kryptograficky overiteľných
identifikátorov.
